\documentclass[aspectratio=169]{beamer}
\usetheme[chinese, code]{Jacquenetta}

\usepackage{amsmath,amssymb}
\usepackage{booktabs}
\usepackage{mwe}   % Provides example-image placeholder for logo

\title{Jacquenetta}
\subtitle{A modern Beamer theme for scientific presentations\\适用于学术演示的现代化极简主题}
\author{Samuel Manchajm}
\institute{GitHub --- Jacquenetta-beamer-theme}
\date{\today}

% Logo placeholder: use an image logo in real documents.
% Text in \logo is evaluated before fonts are ready and may cause nullfont warnings.
\logo{\includegraphics[height=0.7cm]{example-image}}

\begin{document}

% ============================================================
% 标题页 / Title
% ============================================================
\titleframe
\logo{}

% ============================================================
\section{快速开始 / Getting Started}
% ============================================================

\begin{frame}{最小示例 / Minimal example}
    在导言区用一行即可使用 Jacquenetta 主题：

    \vspace{0.1cm}
    {\small\ttfamily
    \textbackslash documentclass[aspectratio=169]\{beamer\}\\
    \textbackslash usetheme\{Jacquenetta\}\quad % 英文内容\\
    \textbackslash usetheme[chinese]\{Jacquenetta\}\quad % 中文内容\\[2pt]
    \textbackslash title\{My Presentation / 我的演示\}\\
    \textbackslash author\{Jane Doe / 张三\}\\
    \textbackslash institute\{University of Somewhere / 某大学\}\\[2pt]
    \textbackslash begin\{document\}\\
    \quad\textbackslash titleframe\\
    \quad\textbackslash begin\{frame\}\{Hello, world! / 你好，世界！\}\\
    \quad\quad Your content here. / 你的内容。\\
    \quad\textbackslash end\{frame\}\\
    \textbackslash end\{document\}
    }
\end{frame}

\begin{frame}{主题选项 / Theme options}
    Jacquenetta 支持以下选项 / Jacquenetta supports the following options:

    \vspace{0.3cm}
    \begin{tabular}{ll}
        \toprule
        \textbf{选项 / Option} & \textbf{效果 / Effect} \\
        \midrule
        \texttt{accent=<颜色 / color>} & 更改强调色 / Change the accent color \\
        \texttt{noborder}       & 关闭签名边框 / Disable the signature border \\
        \bottomrule
    \end{tabular}

    \vspace{0.5cm}
    {\small\ttfamily
    \textbackslash usetheme[accent=teal]\{Jacquenetta\}\\
    \textbackslash usetheme[noborder]\{Jacquenetta\}\\
    \textbackslash usetheme[accent=jqorange, noborder]\{Jacquenetta\}
    }
\end{frame}

% ============================================================
\section{排版 / Typography}
% ============================================================

\sectionframe{01}{排版 Typography}

\begin{frame}{字体样式 / Font styles}
    Jacquenetta 正文使用 \textbf{Helvetica}，中文使用 \textbf{思源黑体}，
    数学公式使用 \textit{衬线字体}。

    \vspace{0.4cm}
    \begin{itemize}
        \item 常规文本 / Regular text for general content
        \item \textbf{粗体文本 / Bold text} for emphasis and headings
        \item \textit{斜体文本 / Italic text} for definitions or foreign terms
        \item \texttt{等宽字体 / Monospace} for code and technical identifiers
        \item \alert{警示文本 / Alerted text} for drawing attention
    \end{itemize}

    \vspace{0.4cm}
    {\footnotesize 脚注大小文本也可用于说明和注释。/ Footnote-sized text is also available for captions and annotations.}
\end{frame}

\begin{frame}{数学公式 / Mathematics}
    Jacquenetta 在数学模式下保留衬线字体以保证可读性。

    \vspace{0.2cm}
    通过贝叶斯定理得到后验概率：
    \[
        P(\theta \mid \mathcal{D}) =
        \frac{P(\mathcal{D} \mid \theta)\, P(\theta)}
             {P(\mathcal{D})}
    \]

    给定 $n$ 个样本，最大似然估计为
    $\hat{\theta} = \arg\max_\theta \mathcal{L}(\theta; x_1, \dots, x_n)$。
\end{frame}

% ============================================================
\section{元素 / Elements}
% ============================================================

\sectionframe{02}{元素 Elements}

\begin{frame}{列表 / Lists}
    \begin{columns}[T]
        \begin{column}{0.45\textwidth}
            \textbf{项目列表 / Itemize}
            \begin{itemize}
                \item 第一层 / First level
                \item 另一项 / Another item
                    \begin{itemize}
                        \item 第二层 / Second level
                        \item 嵌套项 / Nested item
                    \end{itemize}
                \item 回到第一层 / Back to first level
            \end{itemize}
        \end{column}

        \begin{column}{0.45\textwidth}
            \textbf{编号列表 / Enumerate}
            \begin{enumerate}
                \item 第一步 / First step
                \item 第二步 / Second step
                    \begin{enumerate}
                        \item 子步骤 A / Sub-step A
                        \item 子步骤 B / Sub-step B
                    \end{enumerate}
                \item 第三步 / Third step
            \end{enumerate}
        \end{column}
    \end{columns}

    \vspace{0.4cm}
    \textbf{描述列表 / Description}
    \begin{description}
        \item[精确率 / Precision] 预测为正例中真正例的比例。
        \item[召回率 / Recall] 实际正例中被正确预测的比例。
    \end{description}
\end{frame}

\begin{frame}{块环境 / Blocks}
    所有块环境都使用统一的左边框视觉语言，通过颜色区分。

    \begin{block}{默认块 / Default block}
        用于一般信息、定义或中性内容。
    \end{block}

    \begin{alertblock}{警示块 / Alert block}
        用于强调警告、注意事项或重要提示。
    \end{alertblock}

    \begin{exampleblock}{示例块 / Example block}
        用于示例、演示或支持性证据。
    \end{exampleblock}
\end{frame}

\begin{frame}{自定义环境 / Custom environments}
    \begin{problock}{Problock --- 强调色边框标注}
        适用于定理、定义或关键结论。
        对任意事件 $A$ 和 $B$ 且 $P(B) > 0$:\;
        $P(A \mid B) = {P(B \mid A)\, P(A)}/{P(B)}$
    \end{problock}

    \begin{highlightbox}
        \textbf{Highlightbox} --- 橙色高亮标注，用于需要突出的观察、警告或结果。
    \end{highlightbox}
\end{frame}

\begin{frame}{表格 / Tables}
    使用 \texttt{booktabs} 制作简洁专业的表格。

    \vspace{0.4cm}
    \centering
    \begin{tabular}{lccc}
        \toprule
        \textbf{模型} & \textbf{准确率} & \textbf{F1 分数} & \textbf{时间（秒）} \\
        \midrule
        逻辑回归 & 0.89 & 0.87 & 0.3 \\
        随机森林 & 0.93 & 0.91 & 2.1 \\
        神经网络 & 0.95 & 0.94 & 15.7 \\
        \textbf{集成} & \textbf{0.96} & \textbf{0.95} & 18.1 \\
        \bottomrule
    \end{tabular}
\end{frame}

\begin{frame}{分栏 / Columns}
    \begin{columns}[T]
        \begin{column}{0.48\textwidth}
            \begin{problock}{左栏 / Left column}
                分栏适合将文本与图形并排放置，或用于对比两种方法。
            \end{problock}

            \vspace{0.2cm}
            主要优势：
            \begin{itemize}
                \item 更好地利用幻灯片空间
                \item 自然地分隔不同观点
                \item 便于直观对比
            \end{itemize}
        \end{column}

        \begin{column}{0.48\textwidth}
            \begin{problock}{右栏 / Right column}
                你可以混合任何内容：文本、数学、图形、表格或代码列表。
            \end{problock}

            \vspace{0.2cm}
            一个快速公式：
            \[
                \text{ELBO} = \mathbb{E}_{q}[\log p(x \mid z)]
                - D_{\text{KL}}(q \| p)
            \]
        \end{column}
    \end{columns}
\end{frame}

% ============================================================
\section{代码 / Code}
% ============================================================

\sectionframe{03}{代码 Code}

\begin{frame}[fragile]{代码片段 / Code snippet}
    \begin{jqcodebox}[language=Python, title=训练循环 / Training loop]
for epoch in range(num_epochs):
    model.train()
    for batch in train_loader:
        optimizer.zero_grad()
        loss = criterion(model(batch.x), batch.y)
        loss.backward()
        optimizer.step()
    \end{jqcodebox}
\end{frame}

\begin{frame}[fragile]{代码切片 / Code slices}
    \begin{jqlisting}[language=Python, firstline=1, lastline=8, title=Part 1: 导入]
import torch
import torch.nn as nn
from torch.utils.data import DataLoader

class Net(nn.Module):
    def __init__(self):
        super().__init__()
        self.fc = nn.Linear(10, 1)
    \end{jqlisting}
\end{frame}

\begin{frame}[fragile]{代码切片 / Code slices (continued)}
    \begin{jqlisting}[language=Python, firstline=1, lastline=6, firstnumber=9, title=Part 2: 训练]
    def forward(self, x):
        return self.fc(x)

model = Net()
optimizer = torch.optim.Adam(model.parameters(), lr=1e-3)
criterion = nn.MSELoss()
    \end{jqlisting}
\end{frame}

% ============================================================
\section{颜色 / Colors}
% ============================================================

\sectionframe{04}{颜色 Colors}

\begin{frame}{调色板 / Color palette}
    Jacquenetta 定义了四种主题颜色 / Jacquenetta defines four theme colors:

    \vspace{0.4cm}
    \begin{columns}[T]
        \begin{column}{0.22\textwidth}
            \centering
            \begin{tikzpicture}
                \fill[jqdark] (0,0) rectangle (2.5,1.5);
                \node[white] at (1.25,0.75) {\textbf{深色 / Dark}};
            \end{tikzpicture}\\[4pt]
            {\small\texttt{\#121212}}
        \end{column}
        \begin{column}{0.22\textwidth}
            \centering
            \begin{tikzpicture}
                \fill[jqgray] (0,0) rectangle (2.5,1.5);
                \node[white] at (1.25,0.75) {\textbf{灰色 / Gray}};
            \end{tikzpicture}\\[4pt]
            {\small\texttt{\#757575}}
        \end{column}
        \begin{column}{0.22\textwidth}
            \centering
            \begin{tikzpicture}
                \fill[jqblue] (0,0) rectangle (2.5,1.5);
                \node[white] at (1.25,0.75) {\textbf{蓝色 / Blue}};
            \end{tikzpicture}\\[4pt]
            {\small\texttt{\#4A90E2}}
        \end{column}
        \begin{column}{0.22\textwidth}
            \centering
            \begin{tikzpicture}
                \fill[jqorange] (0,0) rectangle (2.5,1.5);
                \node[white] at (1.25,0.75) {\textbf{橙色 / Orange}};
            \end{tikzpicture}\\[4pt]
            {\small\texttt{\#E67E22}}
        \end{column}
    \end{columns}

    \vspace{0.6cm}
    直接使用 / Use them directly: \texttt{\textbackslash color\{jqblue\}} or
    \texttt{\textbackslash textcolor\{jqorange\}\{...\}}。

    \vspace{0.2cm}
    \textbf{强调色 / Accent} color（默认 / default: {\color{jqaccent}blue}）
    控制 / controls the \texttt{problock} border and standard block titles.
\end{frame}

% ============================================================
\section{总结 / Conclusion}
% ============================================================

\sectionframe{05}{总结 Conclusion}

\begin{frame}{总结 / Summary}
    \begin{problock}{Jacquenetta 一览 / Jacquenetta at a glance}
        \begin{itemize}
            \item 极简设计搭配签名深色边框 / Minimalist design with a signature dark border
            \item 干净排版：英文 Helvetica、中文思源黑体、数学衬线 / Clean typography
            \item 自定义环境：\texttt{problock} 和 \texttt{highlightbox}
            \item 使用 \texttt{\textbackslash sectionframe} 的分节页 / Section dividers
            \item 可配置的强调色和可选边框 / Configurable accent color and optional border
        \end{itemize}
    \end{problock}

    \vspace{0.3cm}
    \begin{highlightbox}
        \centering
        \textbf{开始使用 / Get started:}\quad
        \texttt{\textbackslash usetheme[chinese]\{Jacquenetta\}}
    \end{highlightbox}
\end{frame}

\thanksframe[谢谢观看 / Thank you.]

\end{document}
